\subsubsection{RLC Band Stop Filter}

\paragraph{Overview:}
Constructed from a resistor, inductor, and capacitor in series or
parallel, a RLC Band Stop Filter attenuates signals within a
specific frequency range while allowing signals outside this range
to pass. It is a second-order filter.
\paragraph{Schematic:}
\begin{center}
  \begin{circuitikz}
	  % Paths, nodes and wires:
	  \draw (7.5, 5.75) to[american voltage source, l_={$V_{ee}$}] (6.25, 5.75);
	  \node[ground] at (4.75, 6.75){};
	  \node[ground, rotate=-90] at (6.25, 5.75){};
	  \draw (4.75, 8) -| (4.75, 8.25);
	  \draw (8.25, 8.25) -- (9, 8.25);
	  \draw (4.75, 6.75) to[sinusoidal voltage source, l_={$V_{in}$}] (4.75, 8);
	  \node[shape=rectangle, minimum width=0.965cm, minimum height=0.465cm](N1) at (9.5, 8.25){} node[anchor=center] at (N1.text){$V_{out}$};
	  \node[shape=rectangle, minimum width=3.215cm, minimum height=0.465cm] at (6.75, 10.75){} node[anchor=north, align=center, text width=2.827cm, inner sep=6pt] at (6.75, 11){Full Schematic};
	  \node[shape=rectangle, minimum width=4.965cm, minimum height=0.465cm] at (13.5, 10.75){} node[anchor=north, align=center, text width=4.577cm, inner sep=6pt] at (13.5, 11){Small-Signal Schematic};
	  \draw (6, 8.75) to[american inductor, l={$L_1$}] (7.25, 8.75);
	  \draw (6, 7.75) to[capacitor, l_={$C_1$}] (7.25, 7.75);
	  \draw (6, 7.75) -- (5.75, 7.75) -| (5.75, 8.25);
	  \draw (5.75, 8.25) -| (5.75, 8.75) -- (6, 8.75);
	  \draw (7.25, 8.75) -- (7.5, 8.75) -| (7.5, 8.25);
	  \draw (7.25, 7.75) -| (7.5, 8.25);
	  \draw (4.75, 8.25) -- (5.75, 8.25);
	  \draw (8.25, 7.75) to[european resistor, l={$G_1$}] (8.25, 6);
	  \draw (7.5, 5.75) -| (8.25, 6);
	  \draw (8.25, 7.75) -| (8.25, 8.25) -- (7.5, 8.25);
	  \node[ground] at (10.75, 6.75){};
	  \node[ground, rotate=-90] at (12.25, 5.75){};
	  \draw (10.75, 8) -| (10.75, 8.25);
	  \draw (14.25, 8.25) -- (15, 8.25);
	  \draw (10.75, 6.75) to[sinusoidal voltage source, l_={$V_{in}$}] (10.75, 8);
	  \node[shape=rectangle, minimum width=0.965cm, minimum height=0.465cm](N2) at (15.5, 8.25){} node[anchor=center] at (N2.text){$V_{out}$};
	  \draw (12, 8.75) to[american inductor, l={$L_1$}] (13.25, 8.75);
	  \draw (12, 7.75) to[capacitor, l_={$C_1$}] (13.25, 7.75);
	  \draw (12, 7.75) -- (11.75, 7.75) -| (11.75, 8.25);
	  \draw (11.75, 8.25) -| (11.75, 8.75) -- (12, 8.75);
	  \draw (13.25, 8.75) -- (13.5, 8.75) -| (13.5, 8.25);
	  \draw (13.25, 7.75) -| (13.5, 8.25);
	  \draw (10.75, 8.25) -- (11.75, 8.25);
	  \draw (14.25, 7.75) to[european resistor, l={$G_1$}] (14.25, 6);
	  \draw (13.5, 5.75) -| (14.25, 6);
	  \draw (14.25, 7.75) -| (14.25, 8.25) -- (13.5, 8.25);
	  \draw (12.25, 5.75) -- (13.5, 5.75);
	  \draw (7.5, 1.25) to[american voltage source, l_={$V_{ee}$}] (6.25, 1.25);
	  \node[ground] at (4.75, 3){};
	  \node[ground, rotate=-90] at (6.25, 1.25){};
	  \draw (4.75, 4.25) -| (4.75, 4.5);
	  \draw (8.25, 4.5) -- (9, 4.5);
	  \draw (4.75, 3) to[sinusoidal voltage source, l_={$V_{in}$}] (4.75, 4.25);
	  \node[shape=rectangle, minimum width=0.965cm, minimum height=0.465cm](N3) at (9.5, 4.5){} node[anchor=center] at (N3.text){$V_{out}$};
	  \draw (8.25, 4) to[american inductor, l={$L_1$}] (8.25, 2.75);
	  \draw (8.25, 2.75) to[capacitor, l={$C_1$}] (8.25, 1.5);
	  \draw (4.75, 4.5) -- (5.75, 4.5);
	  \draw (7.5, 4.5) to[european resistor, l={$G_1$}] (5.75, 4.5);
	  \draw (7.5, 1.25) -| (8.25, 1.5);
	  \draw (8.25, 4) -| (8.25, 4.5) -- (7.5, 4.5);
	  \node[ground] at (10.75, 3.125){};
	  \node[ground, rotate=-90] at (13.5, 1.375){};
	  \draw (10.75, 4.375) -| (10.75, 4.625);
	  \draw (14.25, 4.625) -- (15, 4.625);
	  \draw (10.75, 3.125) to[sinusoidal voltage source, l_={$V_{in}$}] (10.75, 4.375);
	  \node[shape=rectangle, minimum width=0.965cm, minimum height=0.465cm](N4) at (15.5, 4.625){} node[anchor=center] at (N4.text){$V_{out}$};
	  \draw (14.25, 4.125) to[american inductor, l={$L_1$}] (14.25, 2.875);
	  \draw (14.25, 2.875) to[capacitor, l={$C_1$}] (14.25, 1.625);
	  \draw (10.75, 4.625) -- (11.75, 4.625);
	  \draw (13.5, 4.625) to[european resistor, l={$G_1$}] (11.75, 4.625);
	  \draw (13.5, 1.375) -| (14.25, 1.625);
	  \draw (14.25, 4.125) -| (14.25, 4.625) -- (13.5, 4.625);
	  \draw (8.25, 2.75) -- (9, 2.75);
	  \node[shape=rectangle, minimum width=0.965cm, minimum height=0.465cm](N5) at (9.5, 2.75){} node[anchor=center] at (N5.text){$V_{x}$};
	  \draw (14.25, 2.875) -- (15, 2.875);
	  \node[shape=rectangle, minimum width=0.965cm, minimum height=0.465cm](N6) at (15.5, 2.875){} node[anchor=center] at (N6.text){$V_{x}$};
	  \node[shape=rectangle, minimum width=3.215cm, minimum height=0.465cm] at (2.125, 7.5){} node[anchor=north, align=center, text width=2.827cm, inner sep=6pt] at (2.125, 7.75){Parallel};
	  \node[shape=rectangle, minimum width=3.215cm, minimum height=0.465cm] at (2.25, 3.25){} node[anchor=north, align=center, text width=2.827cm, inner sep=6pt] at (2.25, 3.5){Series};
	  \draw (10, 11) -- (10, 0.5);
	  \draw (1, 5) -- (17, 5);
	  \draw (4, 0.5) -- (4, 11);
	  \draw (1, 10) -- (17, 10);
  \end{circuitikz}
\end{center}

\subsubsubsection{Transfer Function}
\begin{enumerate}
  \item \textbf{Parallel Configuration:}
    Applying KCL at \(V_{out}\) on Full Schematic of Parallel
    Configuration, we get:
    \begin{align*}
      (V_{out} - V_{ee}){G_1} + \frac{(V_{out} - V_{in})}{sL_1} + (V_{out} - V_{in}){sC_1} = 0 \\
      (V_{out} - V_{ee}){G_1} + (V_{out} - V_{in})\left(\frac{1}{sL_1} + sC_1\right) = 0
    \end{align*}
    \paragraph{Superposition:}
      \begin{enumerate}
        \item Due to \(V_{in}\) alone (\(V_{ee} = 0V\)):
          \begin{align*}
            (V_{out} - 0){G_1} + (V_{out} - V_{in})\left(\frac{1}{sL_1} + sC_1\right) &= 0 \\
            V_{out}G_1 + V_{out}\left(\frac{1}{sL_1} + sC_1\right) - V_{in}\left(\frac{1}{sL_1} + sC_1\right) &= 0 \\
            V_{out}\left(G_1 + \frac{1}{sL_1} + sC_1\right) - V_{in}\left(\frac{1}{sL_1} + sC_1\right) &= 0 \\
            V_{out}\left(G_1 + \frac{1}{sL_1} + sC_1\right) &= V_{in}\left(\frac{1}{sL_1} + sC_1\right) \\
            \boxed{H_1(s) = \frac{N_1(s)}{D_1(s)} = \frac{V_{out_1}}{V_{in}} = \frac{\frac{1}{sL_1} + sC_1}{G_1 + \frac{1}{sL_1} + sC_1}}
          \end{align*}
        \item Due to \(V_{ee}\) alone (\(V_{in} = 0V\)):
          \begin{align*}
            (V_{out} - V_{ee}){G_1} + (V_{out} - 0)\left(\frac{1}{sL_1} + sC_1\right) &= 0 \\
            V_{out}G_1 - V_{ee}G_1 + V_{out}\left(\frac{1}{sL_1} + sC_1\right) &= 0 \\
            V_{out}\left(G_1 + \frac{1}{sL_1} + sC_1\right) - V_{ee}G_1 &= 0 \\
            V_{out}\left(G_1 + \frac{1}{sL_1} + sC_1\right) &= V_{ee}G_1 \\
            \boxed{H_2(s) = \frac{N_2(s)}{D_2(s)} = \frac{V_{out_2}}{V_{ee}} = \frac{G_1}{G_1 + \frac{1}{sL_1} + sC_1}}
          \end{align*}
      \end{enumerate}
    \paragraph{Total Parallel Response:}
      \begin{align*}
        V_{out} &= V_{out_1} + V_{out_2} \\
                &= H_1(s)V_{in} + H_2(s)V_{ee} \\
                &= \left(\frac{\frac{1}{sL_1} + sC_1}{G_1 + \frac{1}{sL_1} + sC_1}\right)V_{in} + \left(\frac{G_1}{G_1 + \frac{1}{sL_1} + sC_1}\right)V_{ee} \\
        \boxed{V_{out} = \frac{\left(\frac{1}{sL_1} + sC_1\right)V_{in} + G_1V_{ee}}{G_1 + \frac{1}{sL_1} + sC_1}}
      \end{align*}
  \item \textbf{Series Configuration:}
    Applying KCL at $V_{out}$ on Full Schematic of Series
    Configuration, we get:
    \begin{align*}
      (V_{out} - V_{in})G_1 + \frac{(V_{out} - V_{x})}{sL_1} = 0 \tag{1}
    \end{align*}
    Applying KCL at \(V_{x}\):
    \begin{align*}
      (V_{x} - V_{ee})sC_1 + \frac{(V_{x} - V_{out})}{sL_1} = 0 \tag{2}
    \end{align*}
    From equation (2):
    \begin{align*}
      V_{x}sC_1 - V_{ee}sC_1 + \frac{V_{x}}{sL_1} - \frac{V_{out}}{sL_1} &= 0 \\
      V_{x}\left(sC_1 + \frac{1}{sL_1}\right) &= V_{ee}sC_1 + \frac{V_{out}}{sL_1} \\
      V_{x} &= \frac{V_{ee}sC_1 + \frac{V_{out}}{sL_1}}{sC_1 + \frac{1}{sL_1}} \\
      \boxed{V_{x} = \frac{V_{ee}s^2L_1C_1 + V_{out}}{s^2L_1C_1 + 1}} \tag{3}
    \end{align*}
    \paragraph{Superposition:}
      \begin{enumerate}
        \item Due to \(V_{in}\) alone (\(V_{ee} = 0V\)):
          Substituting equation (3) into equation (1):
          \begin{align*}
            (V_{out} - V_{in})G_1 + \frac{V_{out} - \frac{0 + V_{out}}{s^2L_1C_1 + 1}}{sL_1} &= 0 \\
            (V_{out} - V_{in})G_1 + \frac{V_{out}\left(1 - \frac{1}{s^2L_1C_1 + 1}\right)}{sL_1} &= 0 \\
            (V_{out} - V_{in})G_1 + \frac{V_{out}\left(\frac{s^2L_1C_1}{s^2L_1C_1 + 1}\right)}{sL_1} &= 0 \\
            (V_{out} - V_{in})G_1 + V_{out}\left(\frac{sC_1}{s^2L_1C_1 + 1}\right) &= 0 \\
            V_{out}G_1 - V_{in}G_1 + V_{out}\left(\frac{sC_1}{s^2L_1C_1 + 1}\right) &= 0 \\
            V_{out}\left(G_1 + \frac{sC_1}{s^2L_1C_1 + 1}\right) &= V_{in}G_1 \\
            V_{out}\left(\frac{G_1(s^2L_1C_1 + 1) + sC_1}{s^2L_1C_1 + 1}\right) &= V_{in}G_1 \\
            \boxed{H_3(s) = \frac{N_3(s)}{D_3(s)} = \frac{V_{out_3}}{V_{in}} = \frac{G_1(s^2L_1C_1 + 1)}{G_1(s^2L_1C_1 + 1) + sC_1}}
          \end{align*}
        \item Due to \(V_{ee}\) alone (\(V_{in} = 0V\)):
          Substituting equation (3) into equation (1):
          \begin{align*}
            (V_{out} - 0)G_1 + \frac{V_{out} - \frac{V_{ee}s^2L_1C_1 + V_{out}}{s^2L_1C_1 + 1}}{sL_1} &= 0 \\
            V_{out}G_1 + \frac{V_{out} - \left(\frac{V_{out}}{s^2L_1C_1 + 1} + \frac{V_{ee}s^2L_1C_1}{s^2L_1C_1 + 1}\right)}{sL_1} &= 0 \\
            V_{out}G_1 + \frac{V_{out}\left(1 - \frac{1}{s^2L_1C_1 + 1}\right) - \frac{V_{ee}s^2L_1C_1}{s^2L_1C_1 + 1}}{sL_1} &= 0 \\
            V_{out}G_1 + \frac{V_{out}\left(\frac{s^2L_1C_1}{s^2L_1C_1 + 1}\right) - \frac{V_{ee}s^2L_1C_1}{s^2L_1C_1 + 1}}{sL_1} &= 0 \\
            V_{out}G_1 + \frac{V_{out}}{sL_1}\left(\frac{s^2L_1C_1}{s^2L_1C_1 + 1}\right) - \frac{V_{ee}}{sL_1}\left(\frac{s^2L_1C_1}{s^2L_1C_1 + 1}\right) &= 0 \\
            V_{out}G_1 + V_{out}\left(\frac{sC_1}{s^2L_1C_1 + 1}\right) - V_{ee}\left(\frac{sC_1}{s^2L_1C_1 + 1}\right) &= 0 \\
            V_{out}\left(G_1 + \frac{sC_1}{s^2L_1C_1 + 1}\right) &= V_{ee}\left(\frac{sC_1}{s^2L_1C_1 + 1}\right) \\
            V_{out}\left(\frac{G_1(s^2L_1C_1 + 1) + sC_1}{s^2L_1C_1 + 1}\right) &= V_{ee}\left(\frac{sC_1}{s^2L_1C_1 + 1}\right) \\
            \boxed{H_4(s) = \frac{N_4(s)}{D_4(s)} = \frac{V_{out_4}}{V_{ee}} = \frac{sC_1}{G_1(s^2L_1C_1 + 1) + sC_1}}
          \end{align*}
      \end{enumerate}
    \paragraph{Total Series Response:}
      \begin{align*}
        V_{out} &= V_{out_3} + V_{out_4} \\
                &= H_3(s)V_{in} + H_4(s)V_{ee} \\
                &= \left(\frac{G_1(s^2L_1C_1 + 1)}{G_1(s^2L_1C_1 + 1) + sC_1}\right)V_{in} + \left(\frac{sC_1}{G_1(s^2L_1C_1 + 1) + sC_1}\right)V_{ee} \\
        \boxed{V_{out} = \frac{G_1(s^2L_1C_1 + 1)V_{in} + sC_1V_{ee}}{G_1(s^2L_1C_1 + 1) + sC_1}}
      \end{align*}
\end{enumerate}

\subsubsubsection{Analysis}
Instead of doing a regular analysis, Look at the table below comparing
\(H_1(s), H_2(s), H_3(s), H_4(s)\) of RLC Band Stop Filter (BS) with
\(H_1(s), H_2(s), H_3(s), H_4(s)\) of RLC Band Pass Filter (BP). \\
\begin{tabularx}{\textwidth}{|X|X|X|} \hline
  \textbf{Parameter} & \textbf{BS} & \textbf{BP} \\ \hline
  \(H_1(s)\) & \(\frac{\frac{1}{sL_1} + sC_1}{G_1 + \frac{1}{sL_1} + sC_1}\) & \(\frac{G_1}{G_1 + \frac{1}{sL_1} + sC_1}\) \\ \hline
  \(H_2(s)\) & \(\frac{G_1}{G_1 + \frac{1}{sL_1} + sC_1}\) & \(\frac{\frac{1}{sL_1} + sC_1}{G_1 + \frac{1}{sL_1} + sC_1}\) \\ \hline
  \(H_3(s)\) & \(\frac{G_1(s^2L_1C_1 + 1)}{G_1(s^2L_1C_1 + 1) + sC_1}\) & \(\frac{sC_1}{G_1(s^2L_1C_1 + 1) + sC_1}\) \\ \hline
  \(H_4(s)\) & \(\frac{sC_1}{G_1(s^2L_1C_1 + 1) + sC_1}\) & \(\frac{G_1(s^2L_1C_1 + 1)}{G_1(s^2L_1C_1 + 1) + sC_1}\) \\ \hline
\end{tabularx}\\\\
Do you see the pattern? The transfer functions are simply swapped!
We observe that:
\[
  H_{1,BS}(s) = H_{2,BP}(s), \quad H_{2,BS}(s) = H_{1,BP}(s), \quad H_{3,BS}(s) = H_{4,BP}(s), \quad H_{4,BS}(s) = H_{3,BP}(s)
\]
It also follows for our band stop filter (BS) that, \\\\
\begin{tabularx}{\linewidth}{|X|X|X|X|X|} \hline
  \textbf{Parameter} & \(H_{1,BS}(s)\) & \(H_{2,BS}(s)\) & \(H_{3,BS}(s)\) & \(H_{4,BS}(s)\) \\ \hline
  \textbf{Poles} & Same as \(H_{2,BP}(s)\) & Same as \(H_{1,BP}(s)\) & Same as \(H_{4,BP}(s)\) & Same as \(H_{3,BP}(s)\) \\ \hline
  \textbf{Zeros} & Same as \(H_{2,BP}(s)\) & Same as \(H_{1,BP}(s)\) & Same as \(H_{4,BP}(s)\) & Same as \(H_{3,BP}(s)\) \\ \hline
  \textbf{DC Gain} & Same as \(H_{2,BP}(s)\) & Same as \(H_{1,BP}(s)\) & Same as \(H_{4,BP}(s)\) & Same as \(H_{3,BP}(s)\) \\ \hline
  \textbf{GBP} & Same as Parallel Config & Same as Parallel Config & Same as Series Config & Same as Series Config \\ \hline
  \textbf{Q} & Same as Parallel Config & Same as Parallel Config & Same as Series Config & Same as Series Config \\ \hline
\end{tabularx}\\\\
The impedances are the exact same as in RLC Band Pass Filter.

\subsubsubsection{Question}
Design a RLC Band Stop Filter with a center frequency of
\(f_0 = 1kHz\), bandwidth of \(BW = 100Hz\). Use \(C_1 = 0.1\mu F\).
Compute their input and output impedances at \(f_0\).
\paragraph{Solution:}
\begin{enumerate}
  \item \textbf{Calculate \(L_1\):}
    \begin{align*}
      f_0 &= \frac{1}{2\pi\sqrt{L_1 C_1}} \\
      L_1 &= \frac{1}{(2\pi f_0)^2 C_1} \\
      L_1 &= \frac{1}{(2\pi (1000))^2 (0.1 \times 10^{-6})} \\
      L_1 &\approx 0.2533 H
    \end{align*}
  \item \textbf{Calculate \(R_1\):}
    \begin{align*}
      Q &= \frac{f_0}{BW} = 10 \\
      Q &= R_1 \sqrt{\frac{C_1}{L_1}} \\
      R_1 &= Q \sqrt{\frac{L_1}{C_1}} \\
      R_1 &= 10 \sqrt{\frac{0.2533}{0.1 \times 10^{-6}}} \\
      R_1 &\approx 15915.49 \Omega
    \end{align*}
  \item \textbf{Calculate Input and Output Impedances at \(f_0\):}
    Using the Parallel Configuration:
    \begin{align*}
      Z_{in} = R_1 + \frac{L_1 j \omega}{1 - \omega^2 L_1 C_1}
      \quad \text{and} \quad
      Z_{out} &= \frac{R_1 L_1 j \omega}{R_1(1 - \omega^2 L_1 C_1) + L_1 j \omega}
    \end{align*}
    where \(\omega = 2 \pi f_0 = 6283.19 \, rad/s\).
    \begin{align*}
      Z_{in} &\approx 15915.49 + \frac{0.2533 j (6283.19)}{1 - (6283.19)^2 (0.2533)(0.1 \times 10^{-6})} \\
              &\approx 15915.49 - j 7167.704 \times 10^{15} \Omega \\
      Z_{out} &\approx \frac{15915.49 \times 0.2533 j (6283.19)}{15915.49(1 - (6283.19)^2 (0.2533)(0.1 \times 10^{-6})) + 0.2533 j (6283.19)} \\
              &\approx \frac{25330295.91 j}{-3.533 \times 10^{-12} + j 1591.54} \\
              &\approx \frac{(25330295.91 j)(3.533 \times 10^{-12} - j 1591.54)}{(-3.533 \times 10^{-12} + j 1591.54)(3.533 \times 10^{-12} - j 1591.54)} \\
              &\approx 1.24 \times 10^{-22} - j 15915.5 \Omega
    \end{align*}
\end{enumerate}
